 \documentclass[a4paper, 12pt]{article}
\usepackage[T2A]{fontenc}
\usepackage[utf8]{inputenc}
\usepackage[russian]{babel}
\usepackage{amsmath, amssymb}
\usepackage{physics}
\usepackage{bm}
\usepackage{braket}
\usepackage{hyperref}
\usepackage{graphicx}
\usepackage{xcolor}
\usepackage{subcaption}
\usepackage[verbose=true, a4paper, left=20mm, right=20mm, top=20mm,
bottom=20mm]{geometry}

\title{Конспект лекций по курсу <<Квантовая Механика 1>>  \\ \large Лекция 2}
\author{Лектор: Грибанов Сергей Сергеевич}
\date{8 сентября}

\begin{document}

\maketitle

\tableofcontents

\section{Соотношения неопределенностей}
Поскольку частицы проявляют волновые свойства, то вполне ожидаемо, что для их волновых векторов и частот должны выполняться соотношения неопределенностей по аналогии с соотношениями неопределенностей для электромагнитного излучения, известными из электродинамики. Эти соотношения можно записать в следующем виде
\begin{equation}
  \begin{split}
    \Delta{k}_x\Delta{x} &\gtrsim 1 \Rightarrow \Delta{p}_x\Delta{x} \gtrsim \hslash,\\
    \Delta{k}_y\Delta{y} &\gtrsim 1 \Rightarrow \Delta{p}_y\Delta{y} \gtrsim \hslash,\\
    \Delta{k}_z\Delta{z} &\gtrsim 1 \Rightarrow \Delta{p}_z\Delta{z} \gtrsim \hslash,\\
    \Delta{\omega}\Delta{t} &\gtrsim 1 \Rightarrow \Delta{E}\Delta{t} \gtrsim \hslash,
  \end{split}
\end{equation}
где $\Delta{k}_i$ --- неопределенность в $i$-той компоненте волнового вектора, $\Delta{p}_i$ --- неопределенность в измерении $i$-той компоненты импульса; $\Delta{x}$, $\Delta{y}$ и $\Delta{z}$ --- неопределенности в измерении координат; $\Delta\omega$, $\Delta{E}$ и $\Delta{t}$ --- неопределенности в измерении частоты, энергии и времени. Точный вид соотношений неопределенностей для компонент координаты и импульса
\begin{equation}
  \begin{split}
    \Delta{p}_x\Delta{x} &\geq \hslash / 2,\\
    \Delta{p}_y\Delta{y} &\geq \hslash / 2,\\
    \Delta{p}_z\Delta{z} &\geq \hslash / 2
  \end{split}
\end{equation}
мы получим немного позже.

Таким образом, в отличие от классической механики, в квантовой механике появляются одновременно неизмеримые величины, то есть такие величины, которые нельзя одновременно измерить со сколь угодно большой точностью. Из-за наличия таких величин процессы измерения в квантовой и классической механике существенно отличаются.

Для описания эволюции системы частиц в классической механике достаточно задать значения координат и скоростей частиц в начальный момент времени. Соотношения неопределенности для координаты и импульса утверждают, что мы не можем одновременно сколь угодно точно измерить одноименные компоненты координаты и импульса. Это означает, что в квантовой механике отсутствует понятие траектории. Поскольку понятие траектории отсутствует, то возникает вопрос о том, каким образом описывать состояние квантовомеханической системы и эволюцию этого состояния с течением времени. Ответ на этот вопрос мы дадим чуть позже, введя понятие волновой функции.

\subsection{Оценка наименьшей энергии электрона в атоме водорода}

Атом водорода состоит из протона, вокруг которого движется электрон. Так как
масса протона примерно в $2 \cdot 10^3$ раз больше массы электрона, то протон можно
считать неподвижным. Согласно классической механике, энергия электрона в атоме
водорода имеет вид
\begin{equation}
  E = \frac{\bm{p}^2}{2m_e} - \frac{e^2}{r},
\end{equation}
где $\bm{p}$ --- импульс электрона, $m_e$ --- масса электрона, $e$ --- его заряд,
$r$ --- расстояние от протона до электрона. В классической механике энергия
электрона была бы минимальной в случае падения на центр, когда $\bm{p} = \bm{0}$.
Однако в квантовой механике электрон не может упасть на центр.

Действительно, допустим, электрон движется в малой окрестности $\Delta{r}\equiv{a}$ вблизи
точки $r=0$. Тогда, согласно соотношению неопределенностей
$\Delta{r}\Delta{p}_r\sim\hslash$, неопределенность в импульсе электрона будет не менее чем
$\hslash/a$. Чем меньше $a$, тем большие флуктуации будет испытывать импульс
электрона.

Наименьшая энергия частицы, совершающей финитное движение в некотором поле,
называется энергией основного состояния этой частицы. Оценим энергию основного
состояния электрона в атоме водорода. В этом состоянии электрон должен двигаться
на небольшом расстоянии $r$ от ядра, сопоставимом с неопределенностью этого
расстояния, то есть $r\sim\Delta r = a$.
Аналогично, импульс электрона будет иметь характерное значение порядка
неопределенности импульса, $p \sim \Delta p_r$. Согласно соотношению неопределённостей,
для минимальной неопределенности импульса при заданном $\Delta r = a$ имеем
$\Delta p_r \sim \hslash/a$, откуда $p \sim \hslash/a$. Подставляя эти оценки в выражение для полной
энергии, получим:
\begin{equation}
  \label{eq:hydrogen-atom-energy-as-function-of-a}
  E \sim \frac{\hslash^2}{2m_ea^2} - \frac{e^2}{a}.
\end{equation}

Резкий рост кинетической энергии ($\propto 1/a^2$) при уменьшении $a$ предотвращает
падение электрона на ядро --- фундаментальное следствие соотношения
неопределённостей, отсутствующее в классическом описании.

Найдем, при каком значении параметра $a$
оценка~\eqref{eq:hydrogen-atom-energy-as-function-of-a} энергии электрона
минимальна. Приравняв производную от энергии по параметру $a$ к нулю, получим,
что
\begin{equation}
  -\frac{\hslash^2}{m_ea^3} + \frac{e^2}{a^2} = 0.
\end{equation}
Откуда следует, что значение параметра $a$, отвечающего минимуму энергии
электрона, равно $a_{\text{min}}=\hslash^2/m_ee^2$. Оценка энергии основного
состояния $E_\text{осн}=E(a_{\text{min}})$ принимает следующее значение
\begin{equation}
  \label{eq:hydrogen-E1-estimation}
  E_{\text{осн}} \sim -\frac{m_ee^4}{2\hslash^2}=-\text{Ry}\approx-{13.6}\text{ эВ}.
\end{equation}

Оказывается, что эта оценка совпадает с точным значением энергии основного
состояния атома водорода, которое мы получим позже. Так вышло случайно.
Например, если бы мы написали $\Delta{p}_ra\sim\hslash/2$ вместо $\Delta{p}_ra\sim\hslash$, то
правильный числовой множитель $1/2$ в оценке~\eqref{eq:hydrogen-E1-estimation}
мы бы не получили. Таким образом, использование соотношения неопределенностей
дает оценку энергии основного состояния с точностью до числового множителя.

Параметр $a_{\text{min}}=\hslash^2/m_ee^2$ --- характерный размер атома водорода в
основном состоянии. При более детальном изучении атома водорода мы снова
встретим этот параметр и будем называть его боровским радиусом,
$a_{\text{Б}}=\hslash^2/m_ee^2\approx{0.5\cdot 10^{-8}}\text{ см}$.

\subsection{Соотношение неопределённостей и процесс измерения}

Соотношение неопределённостей проявляется непосредственно в процессе измерения физических величин. Рассмотрим попытку одновременного измерения координаты и импульса электрона. Любое измерение требует взаимодействия с измерительным объектом. В качестве примера можно рассмотреть измерение параметров электрона путём рассеяния на нём фотона.

Чтобы определить координату электрона $x_e$ с точностью $\Delta x_e$, необходимо использовать фотон с длиной волны $\lambda_{\gamma}$ того же порядка величины: $\lambda_{\gamma} \sim \Delta x_e$. Импульс такого фотона определяется соотношением де Бройля:
\[
p_{\gamma} = \frac{2\pi\hslash}{\lambda_{\gamma}} \sim \frac{2\pi\hslash}{\Delta x_e}.
\]
В процессе рассеяния фотона на электроне происходит передача импульса, и согласно закону сохранения импульса, изменение импульса электрона составит величину порядка импульса фотона: $p_e \sim p_{\gamma}$. Следовательно, неопределённость в импульсе электрона после измерения будет:
\[
\Delta p_e \sim p_e \sim \frac{\hslash}{\Delta x_e},
\]
где опущен численный коэффициент $2\pi$ ввиду оценочного характера рассмотрения. Таким образом, чем точнее мы измеряем координату (уменьшая $\Delta x_e$), тем больше становится неопределённость в импульсе электрона.

Если же целью является измерение импульса электрона с малой неопределённостью $\Delta p_e$, следует использовать фотон с большей длиной волны. Поскольку неопределённость в импульсе электрона после взаимодействия удовлетворяет условию $\Delta p_e \sim p_{\gamma} \sim \hslash / \lambda_{\gamma}$, то для уменьшения $\Delta p_e$ необходимо увеличивать $\lambda_{\gamma}$. Однако при этом неопределённость в измерении координаты $\Delta x_e \sim \lambda_{\gamma}$ возрастает. 

Таким образом, существует фундаментальная взаимосвязь между точностью определения координаты и импульса микрочастицы. Невозможно одновременно сколь угодно точно измерить обе эти величины: увеличение точности измерения координаты неизбежно приводит к возрастанию неопределённости импульса, и наоборот.

\section{Дифракция электронов}
Рассмотрим модельный эксперимент по дифракции пучка электронов на щели. Чтобы
дифракционная картина наблюдалась, размер щели должен быть сопоставим с длиной
волны электронов. Схема эксперимента: пучок моноэнергетических электронов
направляется на непрозрачный экран с узкой щелью, за которым расположен
детектор, регистрирующий распределение электронов. В таком эксперименте на
экране наблюдается характерная дифракционная картина с чередованием максимумов и
минимумов интенсивности.

Если бы электроны не проявляли волновых свойств и вели себя исключительно как
классические частицы, то дифракционная картина отсутствовала бы. Вместо
минимумов и максимумов интенсивности наблюдалось бы одно пятно, соответствующее
геометрической проекции щели.

Теперь рассмотрим модификацию этого эксперимента. Будем посылать на щель не
пучок электронов, а отдельные электроны по одному. При прохождении одного
единственного электрона через щель мы не увидим дифракционной картины --- на
экране детектируется лишь одна точка в определенном месте. Однако если повторить
этот эксперимент многократно, регистрируя точки попадания каждого отдельного
электрона, то постепенно на экране проявится та же самая дифракционная картина,
что наблюдалась при использовании пучка электронов.

Этот фундаментальный результат показывает, что дифракционная картина не является
следствием взаимодействия между различными электронами в пучке. Каждый отдельный
электрон проявляет волновые свойства, а наблюдаемое распределение отражает
вероятностный характер квантовых явлений.

Проведем аналогию с электромагнитными волнами. Из курса электродинамики
известно, что интенсивность света в опытах по дифракции пропорциональна квадрату
амплитуды электрического поля: $I(\bm{r}) \propto |\bm{E}(\bm{r})|^2$. В случае электронов мы
также наблюдаем распределение интенсивности, но теперь оно имеет вероятностную
интерпретацию --- плотность точек на экране пропорциональна вероятности
обнаружения электрона в данной области.

Это наводит на мысль ввести для описания состояния частицы некоторую функцию
$\psi(\bm{r}, t)$, аналогичную напряженности электрического поля для электромагнитной
волны. Квадрат модуля этой функции будем интерпретировать как плотность
вероятности обнаружения частицы:
\begin{equation}
  \dv[3]{W}{\bm{r}} = |\psi(\bm{r}, t)|^2
\end{equation}

Эту функцию называют волновой функцией. Для нормированной волновой функции выполняется условие:
\begin{equation}
  \int d^3\bm{r} |\psi(\bm{r}, t)|^2  = 1
\end{equation}
что соответствует требованию, что частица должна быть обнаружена где-либо в
пространстве с единичной вероятностью\footnote{Отметим, что в некоторых случаях интеграл
  от квадрата модуля волновой функции по всему объему расходится, так что
  волновую функцию нельзя нормировать на единицу. Такие случаи мы обсудим
  позже.}.

Таким образом, волновая функция $\psi(\bm{r}, t)$ становится центральным понятием
квантовой механики, описывающим состояние частицы и содержащим информацию о
вероятностях различных результатов измерений.

\section{Интерференция электронов}
Рассмотрим теперь более сложный эксперимент по интерференции пучка электронов,
проходящего через две щели. Данный эксперимент является аналогом классического
опыта Юнга для электромагнитных волн, но проводится с материальными частицами.

Схема эксперимента: пучок электронов направляется на непрозрачный экран с двумя
узкими щелями, за которым расположен детектор. Если закрыть одну из щелей, на
экране наблюдается дифракционная картина от одной щели. При одновременном
открытии обеих щелей возникает система интерференционных полос, которая не
является простой суммой картин от каждой щели в отдельности.

Наблюдение интерференционной картины для пучка электронов, на первый взгляд,
можно было бы попытаться объяснить в рамках корпускулярной модели, предположив,
что она возникает вследствие взаимодействия между разными электронами,
прошедшими через разные щели. Однако этот вывод опровергается модификацией
эксперимента, в которой электроны посылаются на щели поодиночке, с большими
временными интервалами. При каждом акте регистрации на экране детектируется одна
точка, соответствующая попаданию одного электрона. Координата этой точки
случайна. Однако после накопления большого числа таких событий на экране
постепенно проявляется всё та же интерференционная картина. Поскольку электроны
прибывают по одному и не могут взаимодействовать друг с другом, данный результат
однозначно свидетельствует о том, что волновые свойства, приводящие к
интерференции, присущи каждому отдельному электрону.

Данный факт находится в фундаментальном противоречии с классическими
представлениями о траектории частицы. В классической механике частица, проходя
через прибор, должна пройти через одну конкретную щель. Если бы это было так, то
распределение вероятности попадания на экран для одиночных электронов
представляло бы собой сумму распределений от каждой щели, и интерференция
отсутствовала бы. Наблюдение же интерференционной картины для одиночных частиц
означает, что поведение электрона не может быть описано классической
траекторией.

Для дальнейшего анализа рассмотрим попытку экспериментального определения того,
через какую щель проходит электрон. Предположим, что непосредственно за одной из
щелью, например, первой, установлен детектор. При многократном проведении эксперимента с
одиночными электронами обнаруживается, что примерно половина электронов
регистрируется этим детектором, что указывает на их прохождение через первую
щель. Эти электроны поглощаются детектором и не достигают экрана. Оставшаяся
половина электронов, не зарегистрированных детектором, проходит через вторую
щель и формирует на экране дифракционную картину, характерную для одиночной
щели. Интерференционная картина при этом полностью исчезает. Таким образом,
любая попытка определить траекторию частицы, то есть установить, через какую
щель она прошла, разрушает интерференционную картину.

Как было установлено ранее, состояние электрона описывается волновой функцией
$\psi(\bm{r}, t)$, введенной по аналогии с напряженностью электромагнитного поля. Из
электродинамики известно, что электрическое и магнитное поля удовлетворяют
принципу суперпозиции: результирующее поле от нескольких источников представляет
собой сумму полей от каждого источника в отдельности. Естественно предположить,
что волновая функция, также подчиняется принципу суперпозиции.

Эксперименты по наблюдению интерференции различных чатиц (например, электронов)
подтверждают это предположение.
Волновая функция электрона, прошедшего через две щели, должна представлять собой
суперпозицию волновых функций $\psi_1(\bm{r})$ и $\psi_2(\bm{r})$, которые описывали бы
состояние электрона, если бы открыта была только первая или только вторая щель
соответственно:

\begin{equation}
  \psi(\bm{r}) = A_1\psi_1(\bm{r}) + A_2\psi_2(\bm{r})
\end{equation}
где $A_1$ и $A_2$ — комплексные амплитуды, определяемые условиями эксперимента.

Плотность вероятности обнаружения электрона в точке $\bm{r}$ на экране
пропорциональна квадрату модуля волновой функции:
\begin{equation}
  |\psi(\bm{r})|^2 = |A_1\psi_1(\bm{r}) + A_2\psi_2(\bm{r})|^2.
\end{equation}
Раскрывая квадрат модуля, получаем:
\begin{equation}
  |\psi(\bm{r})|^2 = |A_1|^2|\psi_1(\bm{r})|^2 + |A_2|^2|\psi_2(\bm{r})|^2 + 2\Re[A_1A_2^*\psi_1(\bm{r})\psi_2^*(\bm{r})].
\end{equation}
Первые два слагаемых соответствуют распределениям интенсивности, которые
наблюдались бы от каждой щели в отдельности. Последнее слагаемое представляет
собой интерференционный член, ответственный за возникновение интерференционной
картины.

\section{Координатное и импульсное представления}

\subsection{Связь между представлениями}

Ранее, в разделе о дифракции электронов, было введено понятие волновой функции $\psi(\bm{r})$ в координатном представлении, квадрат модуля которой $|\psi(\bm{r})|^2$ задаёт плотность вероятности обнаружения частицы в точке $\bm{r}$. Теперь естественно возникает вопрос: как описать состояние частицы, если нас интересует её импульс? То есть, какова волновая функция $\phi(\bm{p})$ в импульсном представлении, такая что вероятность обнаружить импульс в интервале $[\bm{p}, \bm{p} + \dd[3]{\bm{p}}]$ пропорциональна $|\phi(\bm{p})|^2$? Для установления связи между этими представлениями вернёмся к примеру эксперимента с двумя щелями.

Рассмотрим интерференционную картину от пучка электронов, прошедших через две щели. Волновая функция электрона в точке $\bm{r}$ на экране является суперпозицией двух волн де Бройля, соответствующих прохождению через первую и вторую щели:
\begin{equation}
\psi(\bm{r}) = A_1 e^{i \bm{k}_1 \cdot \bm{r}} + A_2 e^{i \bm{k}_2 \cdot \bm{r}},
\end{equation}
где $\bm{p}_{1,2} = \hslash \bm{k}_{1,2}$ -- импульсы электронов, прошедших через соответствующую щель, а $A_1$ и $A_2$ -- некоторые комплексные амплитуды.

Предположим, мы можем измерить импульс электрона, долетевшего до экрана, с помощью некоего идеального анализатора импульса. Вопрос: какова вероятность $w_1$ зарегистрировать значение импульса $\hslash \bm{k}_1$? Физически это эквивалентно вероятности того, что электрон прилетел именно из первой щели. Опыт и общие принципы квантовой механики подсказывают, что эта вероятность должна быть пропорциональна квадрату модуля амплитуды $A_1$: $w_1 \propto |A_1|^2$. Аналогично, вероятность $w_2$ измерить импульс $\hslash \bm{k}_2$ пропорциональна $|A_2|^2$. Поскольку возможны только эти два исхода, сумма их вероятностей должна быть равна единице. Это условие позволяет определить константу пропорциональности, и в результате вероятности принимают вид:
\begin{align}
w_1 &= \frac{|A_1|^2}{|A_1|^2 + |A_2|^2}, \\
w_2 &= \frac{|A_2|^2}{|A_1|^2 + |A_2|^2}.
\end{align}
Таким образом, амплитуды $A_1$ и $A_2$ можно интерпретировать как амплитуды вероятности обнаружить электрон с импульсом $\hslash \bm{k}_1$ и $\hslash \bm{k}_2$ соответственно.

Действительно, этот результат можно получить и более формально. Нормируем
квадрат модуля волновой функции на единицу в конечном объёме. Поместим нашу
систему в большой параллелепипед с объёмом $V = L_x L_y L_z$, где $L_z$ --
расстояние от щелей до экрана, а $L_x$ и $L_y$ -- поперечные размеры, которые в
дальнейшем мы устремим к бесконечности. Наложим периодические граничные
условия\footnote{Конкретный тип граничных условий вдоль осей $x$ и $y$ не существенен
  для физических результатов, поскольку мы рассматриваем предел $L_x, L_y \to \infty$. В
  таком случае часто бывает удобно использовать периодические граничные
  условия.} на волновую функцию вдоль поперечных координат:
\begin{align*}
\psi(x = -L_x/2, y, z) &= \psi(x = L_x/2, y, z), \\
\psi(x, y = -L_y/2, z) &= \psi(x, y = L_y/2, z).
\end{align*}
Для волновой функции $\psi(\bm{r}) = A_1 e^{i \bm{k}_1 \cdot \bm{r}} + A_2 e^{i \bm{k}_2 \cdot \bm{r}}$ вычислим нормировочный интеграл:
\[
\int_V \dd[3]{\bm{r}}|\psi(\bm{r})|^2  = \int_V \dd[3]{\bm{r}} \left( |A_1|^2 + |A_2|^2 + 2\operatorname{Re}[A_1^* A_2 e^{i (\bm{k}_2 - \bm{k}_1) \cdot \bm{r}}] \right).
\]
Интегралы от осциллирующих членов обращаются в нуль благодаря периодическим граничным условиям по координатам $x$ и $y$, поскольку на интервалах интегрирования $[-L_x/2, L_x/2]$ и $[-L_y/2, L_y/2]$ умещается целое число периодов этих осцилляций. Таким образом, остаётся:
\[
\int_V \dd[3]{\bm{r}} |\psi(\bm{r})|^2 = (|A_1|^2 + |A_2|^2) V.
\]
Потребуем, чтобы полная вероятность была равна единице:
\[
(|A_1|^2 + |A_2|^2) V = 1 \quad \Rightarrow \quad V = \frac{1}{|A_1|^2 + |A_2|^2}.
\]
Теперь вероятность обнаружить электрон с импульсом $\hslash \bm{k}_1$ равна интегралу от квадрата модуля соответствующей волны де Бройля:
\[
w_1 = \int_V \dd[3]{\bm{r}}|A_1 e^{i \bm{k}_1 \cdot \bm{r}}|^2 = |A_1|^2 V = \frac{|A_1|^2}{|A_1|^2 + |A_2|^2}.
\]
Аналогично получаем
\[
w_2 = |A_2|^2 V = \frac{|A_2|^2}{|A_1|^2 + |A_2|^2},
\]
что полностью согласуется с нашими предыдущими выводами.

Обобщим эту идею на случай $N$ щелей. Волновая функция будет суперпозицией $N$ плоских волн:
\begin{equation}
\psi(\bm{r}) = \sum_{n=1}^{N} A_n e^{i \bm{k}_n \cdot \bm{r}}.
\label{eq:n_slits}
\end{equation}
Вероятность обнаружить импульс $\hslash \bm{k}_n$ по-прежнему пропорциональна $|A_n|^2$. Чтобы перейти к непрерывному описанию импульса, представим, что число щелей $N$ становится очень большим, а их расположение таково, что возможные значения импульса $\hslash \bm{k}_n$ заполняют всё пространство импульсов всё более плотно. В пределе $N \to \infty$ дискретная сумма должна превратиться в интеграл по всем возможным волновым векторам $\bm{k}$.

Чтобы совершить этот предельный переход, рассмотрим волновую функцию в большом, но конечном объёме $V = L_x L_y L_z$. На границах этого объёма наложим периодические граничные условия по всем трём направлениям:
\begin{align*}
\psi(x = -L_x/2, y, z) &= \psi(x = L_x/2, y, z), \\
\psi(x, y = -L_y/2, z) &= \psi(x, y = L_y/2, z), \\
\psi(x, y, z = -L_z/2) &= \psi(x, y, z = L_z/2).
\end{align*}
Эти условия приводят к квантованию возможных значений волновых векторов $\bm{k}$:
\begin{align*}
k_x &= \frac{2\pi n_x}{L_x}, \quad n_x = 0, \pm 1, \pm 2, \dots \\
k_y &= \frac{2\pi n_y}{L_y}, \quad n_y = 0, \pm 1, \pm 2, \dots \\
k_z &= \frac{2\pi n_z}{L_z}, \quad n_z = 0, \pm 1, \pm 2, \dots
\end{align*}
Каждое разрешённое значение $\bm{k}$ помечается набором трёх целых чисел $\bm{n} = (n_x, n_y, n_z)$.

Сумму \eqref{eq:n_slits} теперь можно переписать как сумму по всем целым числам $\bm{n}$:
\begin{equation*}
\psi(\bm{r}) = \sum_{\bm{n}} A_{\bm{n}} e^{i \bm{k}_{\bm{n}} \cdot \bm{r}}.
\end{equation*}
Преобразуем эту сумму, введя приращения $\Delta n_x = \Delta n_y = \Delta n_z = 1$:
\begin{equation*}
\sum_{\bm{n}} A_{\bm{n}} e^{i \bm{k}_{\bm{n}} \cdot \bm{r}} = \sum_{\bm{n}} A_{\bm{n}} e^{i \bm{k}_{\bm{n}} \cdot \bm{r}} \Delta n_x \Delta n_y \Delta n_z.
\end{equation*}
Выразим приращения чисел $\Delta n_x$, $\Delta n_y$, $\Delta n_z$ через приращения волновых векторов:
\begin{equation*}
\Delta n_x = \frac{L_x}{2\pi} \Delta k_x, \quad \Delta n_y = \frac{L_y}{2\pi} \Delta k_y, \quad \Delta n_z = \frac{L_z}{2\pi} \Delta k_z.
\end{equation*}
Подставляя эти выражения, получаем:
\begin{equation*}
\sum_{\bm{n}} A_{\bm{n}} e^{i \bm{k}_{\bm{n}} \cdot \bm{r}} \Delta n_x \Delta n_y \Delta n_z = \sum_{\bm{n}} A_{\bm{n}} e^{i \bm{k}_{\bm{n}} \cdot \bm{r}} \frac{L_x L_y L_z}{(2\pi)^3} \Delta k_x \Delta k_y \Delta k_z.
\end{equation*}
В пределе $L_x, L_y, L_z \to \infty$ сумма переходит в интеграл, а $\Delta k_x, \Delta k_y, \Delta k_z$ становятся дифференциалами $\dd k_x, \dd k_y, \dd k_z$. Таким образом, волновую функцию \eqref{eq:n_slits} в этом пределе можно записать как:
\begin{equation}
\psi(\bm{r}) = \int \frac{\dd[3]{\bm{k}}}{(2\pi)^3}\tilde{\phi}(\bm{k}) e^{i \bm{k} \cdot \bm{r}},
\label{eq:fourier_inverse}
\end{equation}
где мы ввели новую функцию:
\begin{equation*}
\tilde{\phi}(\bm{k}) = \lim_{V \to \infty} (A_{\bm{n}} V).
\end{equation*}
Проводя аналогию со случаем двух щелей, где $A_n$ были амплитудами вероятности иметь импульс $\hslash \bm{k}_n$, заключаем, что функция $\tilde{\phi}(\bm{k})$ определяет амплитуды вероятности в пространстве волновых векторов. Выражение \eqref{eq:fourier_inverse} представляет собой \textit{обратное преобразование Фурье}.

Прямое преобразование Фурье, позволяющее выразить $\tilde{\phi}(\bm{k})$ через $\psi(\bm{r})$, имеет вид:
\begin{equation}
\tilde{\phi}(\bm{k}) = \int\dd[3]{\bm{r}} \psi(\bm{r}) e^{-i \bm{k} \cdot \bm{r}}.
\label{eq:fourier_forward}
\end{equation}
Обратите внимание, что в полученных преобразованиях используется
\textit{несимметричное} определение преобразования Фурье: в прямом преобразовании
множитель $(2\pi)^{-3}$ отсутствует, а в обратном -- присутствует. Это означает,
что при интегрировании по волновому вектору $\bm{k}$ в знаменателе всегда будет
присутствовать множитель $(2\pi)^3$, то есть
\begin{equation*}
  \frac{\dd[3]{\bm{k}}}{(2\pi)^3} = \frac{\dd k_x}{2\pi} \frac{\dd k_y}{2\pi} \frac{\dd k_z}{2\pi}.
\end{equation*}
Аналогично, при переходе к интегрированию по импульсу $\bm{p} = \hslash \bm{k}$, учитывая,
что $\dd[3]{\bm{k}} = \dd[3]{\bm{p}} / \hslash^3$, мы будем иметь множитель\footnote{Таким образом,
на каждый дифференциал волнового вектора приходится коэффициент $2\pi$ в
знаменателе, а на каждый дифференциал импульса --- коэффициент $2\pi\hslash$.}
\begin{equation*}
  \frac{\dd[3]{\bm{p}}}{(2\pi\hslash)^3} = \frac{\dd p_x}{2\pi\hslash} \frac{\dd p_y}{2\pi\hslash} \frac{\dd p_z}{2\pi\hslash}.
\end{equation*}
Это соглашение мы будем использовать в дальнейшем.

Зависящие от времени волновые функции преобразуются аналогично. Сначала запишем связь между координатным представлением и представлением в пространстве волновых векторов:
\begin{align}
\psi(\bm{r}, t) &= \int \frac{\dd[3]{\bm{k}}}{(2\pi)^3} \tilde{\phi}(\bm{k}, t) e^{i \bm{k} \cdot \bm{r}}, \label{eq:psi_to_tilde_phi} \\
\tilde{\phi}(\bm{k}, t) &= \int \dd[3]{\bm{r}} \psi(\bm{r}, t) e^{-i \bm{k} \cdot \bm{r}}. \label{eq:tilde_phi_to_psi}
\end{align}
Теперь перейдём от волнового вектора $\bm{k}$ к импульсу $\bm{p}$, вспоминая, что $\bm{k} = \bm{p} / \hslash$. Введём функцию $\phi(\bm{p}, t)$ так, чтобы она совпадала с $\tilde{\phi}(\bm{k}, t)$ при $\bm{k} = \bm{p}/\hslash$, то есть
\begin{equation*}
\phi(\bm{p}, t) = \tilde{\phi}(\bm{k}=\bm{p}/\hslash, t).
\end{equation*}
Тогда, заменяя в преобразованиях \eqref{eq:psi_to_tilde_phi} и \eqref{eq:tilde_phi_to_psi} переменную интегрирования $\bm{k}$ на $\bm{p}/\hslash$, получаем:
\begin{align}
\psi(\bm{r}, t) &= \int \frac{\dd[3]{\bm{p}}}{(2\pi\hslash)^3} \phi(\bm{p}, t) e^{i \bm{p} \cdot \bm{r} / \hslash}, \label{eq:psi_to_phi} \\
\phi(\bm{p}, t) &= \int \dd[3]{\bm{r}} \psi(\bm{r}, t) e^{-i \bm{p} \cdot \bm{r} / \hslash}. \label{eq:phi_to_psi}
\end{align}

Теперь убедимся, что введённая таким образом функция $\phi(\bm{p}, t)$ действительно определяет плотность вероятности обнаружения импульса. Полная вероятность найти частицу где-либо в пространстве равна единице:
\begin{equation}
1 = \int \dd[3]{\bm{r}} |\psi(\bm{r}, t)|^2.
\end{equation}
Подставим сюда выражение для $\psi(\bm{r}, t)$:
\begin{align*}
1 &= \int \dd[3]{\bm{r}} \psi^*(\bm{r}, t) \psi(\bm{r}, t) \\
  &= \int \dd[3]{\bm{r}} \int \frac{\dd[3]{\bm{p}}}{(2\pi\hslash)^3}\phi(\bm{p}, t) \int \frac{\dd[3]{\bm{p}}'}{(2\pi\hslash)^3} \phi^*(\bm{p}', t) e^{i (\bm{p} - \bm{p}') \cdot \bm{r} / \hslash}.
\end{align*}
Поменяем порядок интегрирования, сделав интеграл по $\bm{r}$ внутренним:
\begin{equation*}
1 = \int \frac{\dd[3]{\bm{p}}}{(2\pi\hslash)^3}\phi(\bm{p}, t) \int \frac{\dd[3]{\bm{p}}'}{(2\pi\hslash)^3} \phi^*(\bm{p}', t) \left[ \int \dd[3]{\bm{r}} e^{i (\bm{p} - \bm{p}') \cdot \bm{r} / \hslash} \right].
\end{equation*}
Интеграл в квадратных скобках равен $(2\pi\hslash)^3 \delta^{(3)}(\bm{p} - \bm{p}')$. Используя основное свойство дельта-функции, получаем:
\begin{equation*}
  1 = \int \frac{\dd[3]{\bm{p}}}{(2\pi\hslash)^3} \phi(\bm{p}, t) \int \dd[3]{\bm{p}}' \phi^*(\bm{p}', t) \delta^{(3)}(\bm{p} - \bm{p}') = \int \frac{\dd[3]{\bm{p}}}{(2\pi\hslash)^3} |\phi(\bm{p}, t)|^2
  \equiv \int \dd[3]{\bm{p}} \dv[3]{W(\bm{p}, t)}{\bm{p}}.
\end{equation*}
Таким образом, распределение по импульсам частицы (плотность вероятности в импульсном пространстве) задаётся следующим выражением:
\begin{equation}
\dv[3]{W(\bm{p}, t)}{\bm{p}} = \frac{|\phi(\bm{p}, t)|^2}{(2\pi\hslash)^3}.
\end{equation}

В одномерном случае все формулы упрощаются. Связь между представлениями задаётся преобразованиями:
\begin{align}
\psi(x, t) &= \int \frac{\dd{p}}{2\pi\hslash} \phi(p, t) e^{i p x / \hslash}, \\
\phi(p, t) &= \int \dd{x} \psi(x, t) e^{-i p x / \hslash},
\end{align}
а плотности вероятности определяются как:
\begin{align}
\dv{W(x, t)}{x} &= |\psi(x, t)|^2, \\
\dv{W(p, t)}{p} &= \frac{|\phi(p, t)|^2}{2\pi\hslash}.
\end{align}

\subsection{Операторы в различных представлениях}
В квантовой механике каждой физической величине $A$ ставится в соответствие
линейный оператор $\hat{A}$\footnote{Подробное обсуждение свойств таких операторов будет проведено позже.}. Оператор --- это математический объект, который действует на волновую функцию (состояние), преобразуя её в другую функцию:
\begin{equation}
    \hat{A} \psi(\bm{r}) = \Psi(\bm{r}),
\end{equation}
где $\Psi(\bm{r})$ не обязательно нормирована.

Рассмотрим спектральную задачу для оператора $\hat{A}$, соответствующего физической величине $A$:
\begin{equation}
    \hat{A} \psi_a(\bm{r}) = a \psi_a(\bm{r}).
\end{equation}
Собственные значения $a$ оператора $\hat{A}$ соответствуют значениям, которые может принимать физическая величина $A$ в эксперименте. Если система находится в состоянии, описываемом собственной функцией $\psi_a(\bm{r})$, то измерение величины $A$ даёт определённое значение $a$ (т.е. вероятность получить это значение равна единице). Спектр собственных значений оператора (множество всех возможных $a$) может быть дискретным, непрерывным или смешанным. В случае дискретного спектра собственные значения принимают дискретные (квантованные) значения, и задачу обычно записывают как\footnote{Пока для простоты будем предполагать, что все собственные значения невырождены.}:
\begin{equation}
    \hat{A} \psi_n(\bm{r}) = a_n \psi_n(\bm{r}).
\end{equation}
В случае непрерывного спектра собственное значение может принимать любые значения в некотором интервале.

Рассмотрим действие двух операторов, $\hat{A}$ и $\hat{B}$, на произвольное состояние $\psi(\bm{r})$. Последовательное действие сначала оператора $\hat{A}$, а затем оператора $\hat{B}$ эквивалентно действию некоторого нового оператора $\hat{C}$:
\begin{equation}
    \hat{B} \left( \hat{A} \psi(\bm{r}) \right) = \hat{B} \Psi(\bm{r}) = F(\bm{r}) \equiv \hat{C} \psi(\bm{r}).
\end{equation}
Этот новый оператор называется произведением операторов $\hat{B}$ и $\hat{A}$ и обозначается как $\hat{C} = \hat{B}\hat{A}$.

Одному и тому же оператору $\hat{A}$ в разных представлениях соответствуют разные математические выражения. Однако мы хотим, чтобы его физическое действие, то есть связь между исходным и преобразованным состояниями, было инвариантным относительно выбора представления. Это означает, что если в координатном представлении действие оператора записывается как
\begin{equation}
    \hat{A}_{\bm{r}} \psi(\bm{r}) = \Psi(\bm{r}),
\end{equation}
то в импульсном представлении действие того же самого оператора на соответствующую волновую функцию $\phi(\bm{p})$ должно давать функцию $\Phi(\bm{p})$, которая является Фурье-образом функции $\Psi(\bm{r})$:
\begin{equation}
    \hat{A}_{\bm{p}} \phi(\bm{p}) = \Phi(\bm{p}).
\end{equation}
Для краткости индекс представления у оператора часто опускают, подразумевая, что запись ведётся в том же представлении, что и волновая функция, на которую он действует.

\subsubsection*{Оператор импульса}

Рассмотрим одномерный случай. В импульсном представлении волновая функция $\phi(p)$ является функцией переменной импульса $p$, поэтому оператор импульса в этом представлении просто умножает функцию на $p$:
\begin{equation}
    \hat{p} \phi(p) = p \phi(p), \label{eq:p_in_p_rep}
\end{equation}
то есть $\hat{p} = p$.

Теперь найдём вид этого же оператора в координатном представлении. Для этого воспользуемся связью между представлениями, задаваемой преобразованием Фурье. Пусть $\psi(x)$ --- волновая функция в координатном представлении, а $\phi(p)$ --- её Фурье-образ:
\begin{equation}
    \phi(p) = \int_{-\infty}^{\infty} \dd{x} \psi(x) e^{-i p x / \hslash}.
\end{equation}
Обратное преобразование Фурье имеет вид:
\begin{equation}
    \psi(x) = \int_{-\infty}^{\infty} \frac{\dd{p}}{2\pi\hslash} \phi(p) e^{i p x / \hslash}.
\end{equation}

Действие оператора импульса на волновую функцию в координатном представлении можно выразить через его действие в импульсном представлении:
\begin{align}
    \hat{p} \psi(x) &= \int_{-\infty}^{\infty} \frac{\dd{p}}{2\pi\hslash} \left[\hat{p} \phi(p)\right] e^{i p x / \hslash} \nonumber \\
                    &= \int_{-\infty}^{\infty} \frac{\dd{p}}{2\pi\hslash} \left[p \phi(p)\right] e^{i p x / \hslash} \nonumber \\
                    &= \int_{-\infty}^{\infty} \frac{\dd{p}}{2\pi\hslash} \phi(p) \left[p e^{i p x / \hslash}\right] \nonumber \\
                    &= \int_{-\infty}^{\infty} \frac{\dd{p}}{2\pi\hslash} \phi(p) \left[-i\hslash \pdv{x} e^{i p x / \hslash}\right] \nonumber \\
                    &= -i\hslash \pdv{x} \int_{-\infty}^{\infty} \frac{\dd{p}}{2\pi\hslash} \phi(p) e^{i p x / \hslash} \nonumber \\
                    &= -i\hslash \pdv{x} \psi(x).
\end{align}
Таким образом, оператор импульса в координатном представлении имеет вид:
\begin{equation}
    \hat{p} = -i \hslash \pdv{x}. \label{eq:p_in_x_rep}
\end{equation}

\subsubsection*{Оператор координаты}

Аналогичные рассуждения можно провести для оператора координаты $\hat{x}$. В координатном представлении его действие сводится к умножению на координату $x$:
\begin{equation}
    \hat{x} \psi(x) = x \psi(x), \label{eq:x_in_x_rep}
\end{equation}
то есть $\hat{x} = x$.

Чтобы найти вид этого оператора в импульсном представлении, выразим его действие через прямое преобразование Фурье:
\begin{align}
    \hat{x} \phi(p) &= \int_{-\infty}^{\infty} \dd{x} \left[\hat{x} \psi(x)\right] e^{-i p x / \hslash} \nonumber \\
                    &= \int_{-\infty}^{\infty} \dd{x} \left[x \psi(x)\right] e^{-i p x / \hslash} \nonumber \\
                    &= \int_{-\infty}^{\infty} \dd{x} \psi(x) \left[x e^{-i p x / \hslash}\right] \nonumber \\
                    &= \int_{-\infty}^{\infty} \dd{x} \psi(x) \left[i\hslash \pdv{p} e^{-i p x / \hslash}\right] \nonumber \\
                    &= i\hslash \pdv{p} \int_{-\infty}^{\infty} \dd{x} \psi(x) e^{-i p x / \hslash} \nonumber \\
                    &= i\hslash \pdv{p} \phi(p).
\end{align}
Таким образом, оператор координаты в импульсном представлении имеет вид:
\begin{equation}
    \hat{x} = i \hslash \pdv{p}. \label{eq:x_in_p_rep}
\end{equation}

\subsubsection*{Обобщение на трёхмерный случай}

Полученные результаты обобщаются на трёхмерный случай следующим образом:

\begin{center}
\begin{tabular}{lll}
\hline
Представление & Оператор координаты & Оператор импульса \\
\hline
Координатное & $\hat{\bm{r}} = \bm{r}$ & $\hat{\bm{p}} = -i\hslash \bm{\nabla}$ \\
Импульсное & $\hat{\bm{r}} = i\hslash \pdv*{\bm{p}}$ & $\hat{\bm{p}} = \bm{p}$ \\
\hline
\end{tabular}
\end{center}

где $\bm{\nabla} = \left( \pdv{x}, \pdv{y}, \pdv{z} \right)$ --- градиент, а
$\pdv{\bm{p}} = \left( \pdv{p_x}, \pdv{p_y}, \pdv{p_z} \right)$ --- векторный оператор
дифференцирования по компонентам импульса.

\section{Уравнение Шредингера}

Для описания пространственно-временной эволюции волн де Бройля необходимо динамическое уравнение, определяющее изменение волновой функции в пространстве и времени. Мы приведем физическое обоснование вида такого уравнения, используя представление о волновой природе квантовых объектов.

Рассмотрим волновой процесс, описываемый комплекснозначной функцией
$\psi(\bm{r}, t)$. В классической физике распространение волн описывается волновым
уравнением\footnote{Здесь и далее оператор $\nabla^2$ обозначает лапласиан (оператор Лапласа),
  который в декартовой системе координат имеет вид $\nabla^2 = \Delta = \pdv[2]{x} + \pdv[2]{y} + \pdv[2]{z}$.}:
\[
\nabla^2 \psi - \frac{1}{v^2_{\phi}} \pdv[2]{\psi}{t} = 0,
\]
где $v_{\phi}$ -- фазовая скорость волны.

Фазовая скорость волны де Бройля, связанной с частицей массы $m$, выражается через ее энергию и импульс. Используя соотношения де Бройля, получаем:
\[
v_{\phi} = \frac{\omega}{k} = \frac{\hslash\omega}{p} = \frac{E}{p}.
\]
Для нерелятивистской частицы кинетическая энергия связана с импульсом соотношением:
\[
K = \frac{p^2}{2m} = E - U(\bm{r}),
\]
поэтому фазовая скорость равна:
\[
v_{\phi} = \frac{E}{\sqrt{2m(E - U(\bm{r}))}}.
\]

Заметим, что фазовая скорость зависит от координат через потенциальную энергию $U(\bm{r})$, что аналогично распространению волн в среде с переменным показателем преломления.

Рассмотрим теперь частный случай, когда волновая функция представляет собой суперпозицию волн де Бройля с одной и той же частотой. Такой волновой пакет можно записать в виде:
\[
\psi_E(\bm{r}, t) = \int \frac{\dd[3]{\bm{p}}}{(2\pi\hslash)^3} \phi_E(\bm{p}) e^{i(\bm{p}\cdot\bm{r} - Et)/\hslash} = \psi_E(\bm{r}) e^{-iEt/\hslash},
\]
где $E = \hslash\omega$ -- определенное значение энергии. В этом случае плотность вероятности $|\psi_E(\bm{r}, t)|^2 = |\psi_E(\bm{r})|^2$ не зависит от времени. Состояния с такой временной зависимостью называются \emph{стационарными}.

Из вида волновой функции стационарного состояния следует важное соотношение:
\begin{equation}
E\psi_E = i\hslash \pdv{\psi_E}{t}.
\end{equation}

Подставляя выражение для фазовой скорости в волновое уравнение и используя вид волновой функции для стационарных состояний, приходим к стационарному уравнению Шредингера:
\begin{equation}
\left( -\frac{\hslash^2}{2m} \nabla^2 + U(\bm{r}) \right) \psi_E(\bm{r}) = E \psi_E(\bm{r}).
\end{equation}

Для общего случая волновой функции произвольного нестационарного состояния имеет место нестационарное уравнение Шредингера:
\begin{equation}
i\hslash \pdv{\psi(\bm{r}, t)}{t} = \left( -\frac{\hslash^2}{2m} \nabla^2 + U(\bm{r}, t) \right) \psi(\bm{r}, t).
\end{equation}

Заметим, что выражение $-\hslash^2\nabla^2$ представляет собой квадрат оператора импульса $\hat{\bm{p}}^2$, поскольку в координатном представлении $\hat{\bm{p}} = -i\hslash\bm{\nabla}$. Это позволяет ввести \emph{оператор Гамильтона} (гамильтониан):
\begin{equation}
\hat{H} = \frac{\hat{\bm{p}}^2}{2m} + U(\bm{r}, t) = -\frac{\hslash^2}{2m} \nabla^2 + U(\bm{r}, t).
\end{equation}

С использованием гамильтониана уравнения Шредингера принимают компактный вид. Нестационарное уравнение Шредингера записывается как:
\begin{equation}
i\hslash \pdv{\psi(\bm{r}, t)}{t} = \hat{H} \psi(\bm{r}, t),
\end{equation}
а стационарное уравнение Шредингера представляет собой уравнение на собственные функции и собственные значения оператора Гамильтона:
\begin{equation}
\hat{H} \psi_E(\bm{r}) = E \psi_E(\bm{r}).
\end{equation}
Собственные значения $E$ соответствуют возможным значениям энергии системы, а собственные функции $\psi_E(\bm{r})$ описывают стационарные состояния с определенной энергией.

Отметим фундаментальные свойства уравнения Шредингера:

\begin{enumerate}
\item Уравнение Шредингера линейно относительно волновой функции. Это означает, что любая линейная комбинация решений этого уравнения с произвольными комплексными коэффициентами также является его решением. Таким образом, волновые функции удовлетворяют принципу суперпозиции.

\item Поскольку уравнение Шредингера содержит первую производную по времени, эволюция системы полностью определяется заданием начального условия. Если известна волновая функция $\psi(\bm{r}, t_0)$ в некоторый начальный момент времени $t_0$, то уравнение однозначно определяет волновую функцию $\psi(\bm{r}, t)$ в любой последующий момент времени $t > t_0$.
\end{enumerate}


\subsection{Стационарные состояния}

Рассмотрим более формально свойства стационарных состояний и их зависимость от времени. По определению, волновая функция стационарного состояния удовлетворяет стационарному уравнению Шредингера:
\begin{equation}
\hat{H}\psi_E = E \psi_E.
\end{equation}
С другой стороны, как и любая волновая функция, она должна быть решением нестационарного уравнения Шредингера:
\begin{equation}
\hat{H}\psi_E = i\hslash \pdv{\psi_E}{t}.
\end{equation}
Приравнивая правые части этих уравнений, получаем:
\begin{equation}
i\hslash \pdv{\psi_E}{t} = E \psi_E.
\end{equation}
Это дифференциальное уравнение легко решается и даёт искомую зависимость от времени. Учитывая произвол в выборе начала отсчёта времени, положим для простоты $t_0 = 0$, тогда решение имеет вид:
\begin{equation}
\psi_E(\bm{r}, t) = \psi_E(\bm{r}, 0) e^{-i E t / \hslash} = \psi_E(\bm{r}) e^{-i E t / \hslash},
\end{equation}
где $\psi_E(\bm{r}) \equiv \psi_E(\bm{r}, 0)$ --- пространственная часть волновой функции.

Квадрат модуля волновой функции стационарного состояния не зависит от времени:
\begin{equation}
|\psi_E(\bm{r}, t)|^2 = |\psi_E(\bm{r}) e^{-i E t / \hslash}|^2 = |\psi_E(\bm{r})|^2,
\end{equation}
поскольку комплексная экспонента с мнимым показателем по модулю равна единице. Это означает, что плотность вероятности нахождения частицы в различных точках пространства остаётся постоянной во времени.

\subsection{Эволюция произвольного состояния}

Зависимость от времени волновой функции в общем случае определяется нестационарным уравнением Шредингера. Для простоты предположим, что спектр энергий дискретный, так что собственные значения оператора Гамильтона можно пронумеровать: $\hat{H}\psi_n(\bm{r}) = E_n \psi_n(\bm{r})$.

Позднее будет показано, что собственные функции оператора Гамильтона образуют базис в гильбертовом пространстве. Это означает, что произвольную волновую функцию в начальный момент времени $t=0$ можно разложить в ряд по этим функциям:\footnote{В случае непрерывного спектра суммирование заменяется на интегрирование по энергии. С этим случаем мы познакомимся подробнее позднее.}
\begin{equation}
\psi(\bm{r}, 0) = \sum_n c_n(0) \, \psi_n(\bm{r}),
\label{eq:initial_decomposition}
\end{equation}
где коэффициенты $c_n(0)$ определяются видом начальной волновой функции $\psi(\bm{r}, 0)$.

В произвольный момент времени решение ищется в виде аналогичного разложения, но с коэффициентами, зависящими от времени:
\begin{equation}
\psi(\bm{r}, t) = \sum_n c_n(t) \, \psi_n(\bm{r}).
\label{eq:time_decomposition}
\end{equation}
Подставим это разложение в нестационарное уравнение Шредингера $i\hslash \pdv*{\psi}{t} = \hat{H}\psi$:
\begin{equation}
i\hslash \sum_n \dot{c}_n(t) \psi_n(\bm{r}) = \sum_n c_n(t) \, \hat{H}\psi_n(\bm{r}) = \sum_n c_n(t) E_n \psi_n(\bm{r}).
\end{equation}
Перенося все члены в левую часть, получаем:
\begin{equation}
\sum_n \left[ i\hslash \dot{c}_n(t) - E_n c_n(t) \right] \psi_n(\bm{r}) = 0.
\end{equation}
Поскольку функции $\psi_n(\bm{r})$ линейно независимы (образуют базис), это равенство выполняется только тогда, когда коэффициент перед каждой функцией равен нулю:
\begin{equation}
i\hslash \dot{c}_n(t) = E_n c_n(t) \quad \text{для всех } n.
\end{equation}
Решение этого уравнения с начальным условием $c_n(0)$ из разложения \eqref{eq:initial_decomposition} имеет вид:
\begin{equation}
c_n(t) = c_n(0) \, e^{-i E_n t / \hslash}.
\end{equation}
Подставляя это решение в разложение \eqref{eq:time_decomposition}, окончательно получаем закон эволюции произвольного состояния:
\begin{equation}
\psi(\bm{r}, t) = \sum_n c_n(0) \, e^{-i E_n t / \hslash} \, \psi_n(\bm{r}) = \sum_n c_n(0) \, \psi_n(\bm{r}, t),
\end{equation}
где $\psi_n(\bm{r}, t) = \psi_n(\bm{r}) \, e^{-i E_n t / \hslash}$ --- это зависимость от времени волновой функции $n$-го стационарного состояния.

Таким образом, временная эволюция волновой функции полностью определяется её
видом в начальный момент времени, что соответствует фундаментальному
свойству уравнения Шредингера как дифференциального уравнения первого порядка по
времени.


\end{document}
